\par\noindent\SourceLeaf{46}%
und das allgemeine Glied $r^{\text{tus}}$

\[ = \frac{1}{2^{2m-1}}{}^{2m}\!\mathfrak{B}^{(m+r)}\! . \cos 2r\varphi. \]

$\beta$. Die ungeraden Potenzen folgen, wie sich zeigt, demselben Gesetz, außer dass
\noindent der Zähler des Anfangsgliedes den ganzen späteren der mittleren Binomialkoeffizienten der gegebenen Potenz enthält (bei ungeradem Exponenten treten nämlich zwei mittlere und einander gleiche Koeffizienten auf). Das allgemeine Gesetz der Reihe lautet also wie folgt. Die Reihe für $\cos\varphi^{2m-1}$ hat das
\noindent Anfangsglied

\[ = \frac{1}{2^{2m-2}}{}^{2m-1}\!\mathfrak{B}^{(m)}\cos\varphi, \]

\noindent und das allgemeine Glied $r^{\text{tus}}$

\[ = \frac{1}{2^{2m-2}}{}^{2m-1}\!\mathfrak{B}^{(m+r)}\cos(2r+1)\varphi. \]

3. Nehmen wir nun zunächst an, dass das vorgeschlagene Gesetz für eine beliebige ungerade Potenz gilt, z.\,B.\ für die Potenz $(2m-1)^{\text{ta}}$, und dass somit gilt

\[
\begin{aligned}
\cos\varphi^{2m-1} & = \frac{1}{2^{2m-2}}{}^{2m-1}\!\mathfrak{B}^{(m)}\cos\varphi +\\
&{} \frac{1}{2^{2m-2}}{}^{2m-1}\!\mathfrak{B}^{(m+1)}\cos 3\varphi +\\
&{} \frac{1}{2^{2m-2}}{}^{2m-1}\!\mathfrak{B}^{(m+2)}\cos 5\varphi + \text{etc.},
\end{aligned}
\]

\noindent Die um eins höhere gerade Potenz erhält man, indem man beide Seiten der Gleichung mit $\cos\varphi$ multipliziert. Jedes auf diese Weise gebildete Glied der Reihe enthält ein Produkt aus zwei Kosinus, das sich nach dem bekannten Satz leicht in die Summe zweier Kosinus auflösen lässt. So ergibt sich

\[
\begin{aligned}
\cos\varphi^{2m} & = \frac{1}{2^{2m-2}}{}^{2m-1}\!\mathfrak{B}^{(m)}\cos\varphi . \cos\varphi +\\
&{} \frac{1}{2^{2m-2}}{}^{2m-1}\!\mathfrak{B}^{(m+1)}\cos 3\varphi . \cos\varphi +\\
&{} \frac{1}{2^{2m-2}}{}^{2m-1}\!\mathfrak{B}^{(m+2)}\cos 5\varphi . \cos\varphi + .\,.\,.
\end{aligned}
\]

\noindent Da nun

\[ \cos\varphi\cos\varphi = \tfrac12+\tfrac12\cos 2\varphi, \]

\noindent wird das Anfangsglied zu

\[ = \frac{1}{2^{2m-2}}{}^{2m-1}\!\mathfrak{B}^{(m)} . \tfrac12 +
\frac{1}{2^{2m-2}}{}^{2m-1}\!\mathfrak{B}^{(m)}\frac{\cos 2\varphi}{2}, \]

\noindent und nach demselben Gesetz wird das erste Glied zu

\[ = \frac{1}{2^{2m-2}}{}^{2m-1}\!\mathfrak{B}^{(m+1)} . \frac{\cos 2\varphi}{2} +
\frac{1}{2^{2m-2}}{}^{2m-1}\!\mathfrak{B}^{(m+1)}\frac{\cos 4\varphi}{2}. \]
